% A.66: arithmetic input for the fixed-depth lifting theorem A.65.
\begin{statement}{A.66}{Proved}{Ordinary infinite-support pairs with any fixed product of tame denominators}
Fix integers $d\ge2$ and $C>0$ with $24\mid C$.
Let $J\ge2d^2$ be at least every prime divisor of $C$, and put
\[
 n_j=(j+J)!,\qquad D_j=Cd!n_j^d,\qquad
 E_j(Q)=Cn_j\prod_{r=2}^{d}(Q^r-1).
 \tag{A.66.1}
\]
There are three sequences of distinct primes $q_j,r_j,s_j$, all
greater than three and with no repetitions between the sequences,
for which the completed classes
\[
 A=3\prod_{j\ge0}\frac{q_j^{E_j(q_j)}}{s_j^{E_j(s_j)}},
 \qquad
 B=-2\prod_{j\ge0}\frac{r_j^{E_j(r_j)}}{s_j^{E_j(s_j)}}
 \tag{A.66.2}
\]
form an infinite-support ordinary completed unit-equation pair.
At every finite prime $p$ there is $x_p\in\Q_p\setminus\{0,1\}$
whose two Kummer classes are the restrictions of $A,B$.
At the real place the point $3$ has these two classes.

The valuation patterns at two and three are $(0,1)$ and $(1,0)$.
At the three new primes of stage $j$ they are respectively
\[
 (E_j(q_j),0),\qquad(0,E_j(r_j)),\qquad
 (-E_j(s_j),-E_j(s_j));
 \tag{A.66.3}
\]
all other valuations vanish.  In particular the pair comes from
neither a rational point nor a cusp packet.  There are explicit
completed classes $A_0,B_0$ with $A=3A_0^C$ and
$B=-2B_0^C$.  Its image in
$(\Q^\times/(\Q^\times)^C)^2$ is exactly $(3,-2)$.
Thus $C$ can impose any prescribed finite abelian congruence to
the point $3$ while retaining the displayed denominator factors.
\end{statement}
\begin{proof}
\emph{The exact local image at a prime dividing the factorial.}
Write $n=n_j$ and fix $p\mid n$.  Set
\[
 s=v_p(n),\quad a=v_p(d!),\quad b=v_p(d),\quad
 A_p=Cnd!p^{ds},\quad h=v_p(A_p)=(d+1)s+v_p(Cd!).
 \tag{A.66.4}
\]
On $z=1+p^s w$ consider the convergent function
\[
 F(z)=\exp\left(Cn\prod_{r=2}^d(z^r-1)\log z\right).
 \tag{A.66.5}
\]
The inequalities below give $s\ge4$ at two, so both logarithm
and exponential converge on all the indicated balls.
Define
\[
 P(t)=\prod_{r=2}^d\frac{(1+t)^r-1}{rt},\qquad
 V(t)=P(t)\frac{\log(1+t)}t.
 \tag{A.66.6}
\]
The polynomial $d!P(t)$ has integral coefficients and
$P(0)=V(0)=1$.  The coefficient of $t^k$ in $V$ has valuation
at least $-a-\lfloor\log_p(k+1)\rfloor$: each summand is a
coefficient of $P$ multiplied by one of
$1,1/2,\ldots,1/(k+1)$.  Consequently
\[
 V(p^sw)\in1+p^\kappa\Z_p\langle w\rangle,
 \qquad
 \kappa=
 \begin{cases}s-a-1,&p\le d,\\s,&p>d.\end{cases}
 \tag{A.66.7}
\]
Here $\Z_p\langle w\rangle$ denotes restricted power series.
For the first bound use
$\lfloor\log_p(k+1)\rfloor\le k$ and $s\ge1$.
For the second, $a=0$, $p\ge3$, and
$\lfloor\log_p(k+1)\rfloor\le k-1$ for $k\ge1$.
The actual coefficient bounds tend to infinity with $k$.

For $p\le d$ the explicit factorial choice gives
\[
 s\ge v_p(J!)\ge\lfloor J/p\rfloor\ge2d
           \ge v_p(d!)+v_p(d)+2=a+b+2.
 \tag{A.66.8}
\]
Indeed $a,b\le d-1$.  For $p>d$ we have $a=b=0$ and
$s\ge1$.  In every case $\kappa\ge b+1$.
Expanding (A.66.5) gives
\[
 F(1+p^s w)=1+A_p f(w),\qquad
                 f(w)=w^d+p^\kappa R(w),
            \quad R\in\Z_p\langle w\rangle.
 \tag{A.66.9}
\]
To verify the exponential terms, the coefficient factor of
exponential degree $u\ge2$, after division by $A_p$, has
valuation $(u-1)h-v_p(u!)\ge(u-1)(h-1)\ge\kappa$.
Thus no unexamined logarithmic or exponential denominator
enters (A.66.9).

Put $z_0=1+p^s$.  On $w=1+p^{b+1}t$ the normalized map is
\[
 \frac{f(1+p^{b+1}t)-f(1)}{p^{2b+1}}
                 =\frac{d}{p^b}t+pG(t),
                    \qquad G\in\Z_p\langle t\rangle.
 \tag{A.66.10}
\]
For every binomial term of degree $k\ge2$ its coefficient
valuation after division is at least
$(b+1)(k-1)-b\ge1$; the error term in (A.66.9) contributes
at least $\kappa-b\ge1$.
A restricted series $ut+pG(t)$ with $u\in\Z_p^\times$ is
bijective on $\Z_p$: for a prescribed value, iteration of
$t\mapsto u^{-1}(y-pG(t))$ is a strict contraction.
We therefore have the exact image and derivative formulas
\[
 \begin{split}
 F\bigl(z_0+p^{s+b+1}\Z_p\bigr)
       &=F(z_0)+p^{h+2b+1}\Z_p,\\
 v_p(F'(z_0))&=\delta_j(p):=ds+v_p(Cd!)+b.
 \end{split}
 \tag{A.66.11}
\]
The derivative identity uses $f'(1)=d+p^\kappa R'(1)$
and $\kappa>b$.  The entire input ball has $v_p(z-1)=s$.

\emph{The all-place invariant.}
For $p\nmid n_j$ set $\delta_j(p)=0$, and for every prime put
\[
 H_j(p)=2\delta_j(p)+j+3,\qquad B(D)=(D+1)^4.
 \tag{A.66.12}
\]
We construct integers
\[
 a_j=3u_j^{24}>0,\qquad b_j=-2v_j^{24}<0,
                    \qquad c_j=w_j^{24}>0,
 \tag{A.66.13}
\]
with pairwise coprime supports and the following properties:
$a_jX^{D_j}+b_jY^{D_j}=c_jZ^{D_j}$ has a solution
in $(\Z_p^\times)^3$ at every prime $p$, and
\[
 v_p(a_j+b_j-c_j)\ge H_j(p)
                              \quad(p\le B(D_j)).
 \tag{A.66.14}
\]
Start from $(a_0,b_0,c_0)=(3,-2,1)$.
At a prime dividing $n_j$, the difference between $H_j(p)$
and the image depth in (A.66.11) is
\[
 H_j(p)-(h+2b+1)
            =(d-1)s+v_p(Cd!)+j+2>0.
 \tag{A.66.15}
\]
Also $H_j(p)>2v_p(D_j)$; the same latter inequality holds
at $p\nmid n_j$, since every prime divisor of $Cd!$
already divides $n_j$.

\emph{The finite triangular system of congruences.}
At stage $j$ choose an integer $P_0$ at least $B(D_{j+1})$
and at least every prime dividing $a_jb_jc_j$.  Process
every prime $p\le P_0$ in increasing order.
For each of the three future primes $Q$, impose
\[
 \begin{cases}
 v_p(Q-1)=v_p(n_j),&p\mid n_j,\\
 Q\ne0,\quad Q^r\ne1\pmod p\quad(2\le r\le d),
                                               &p\nmid n_j.
 \end{cases}
 \tag{A.66.16}
\]
Choose the modulus $p^{L_p}$ large enough that
\[
 L_p\ge H_{j+1}(p),\qquad p^{L_p}>P_0,
 \qquad L_p\ge v_p(n_j)+v_p(d)+2\quad(p\mid n_j).
 \tag{A.66.17}
\]
The second condition ensures that, before processing $p$,
the residues modulo every prime-power divisor of $p-1$
have already been fixed.  Hence the residue
$e_i=E_j(Q_i)\pmod{p-1}$ is already fixed separately for
each of the three future primes.

At a prime $\ell\mid n_j$, the elementary lifting-the-exponent
identity gives
\[
 v_\ell\left(\prod_{r=2}^d(Q_i^r-1)\right)
       =(d-1)v_\ell(n_j)+v_\ell(d!),
 \qquad v_\ell(E_j(Q_i))=v_\ell(D_j).
 \tag{A.66.18}
\]
At two this uses $Q_i\equiv1\pmod{16}$, which follows
from (A.66.8).  At every already processed prime
$\ell\nmid n_j$, the exclusions in (A.66.16) make both
valuations zero.  Therefore
\[
                  \gcd(e_i,p-1)=\gcd(D_j,p-1).
 \tag{A.66.19}
\]
The three exponent maps on $\F_p^\times$ have precisely
the same image, the subgroup of $D_j$th powers.

\emph{Choosing the local residues and their arbitrary precision.}
At $p\mid n_j$ use (A.66.11) with the same central input
$z_0$ in all three coordinates.  At least one of
$a_j,b_j,c_j$ is a unit.  Varying its input realizes
every correction in $p^{h+2b+1}\Z_p$, after multiplication
by that unit coefficient.  Equations (A.66.14)--(A.66.15)
therefore give three permitted local inputs satisfying
\[
                   a_jF(q)+b_jF(r)=c_jF(s).
 \tag{A.66.20}
\]

For $p\nmid n_j$, every primitive residue has order
$p-1>d$, because $p>j+J\ge2d^2$.
Fix a noncritical residue $\omega_i\in\mu_{p-1}\subset\Z_p$
for each coordinate.  On its residue ball use the exact map
\[
 F_{i,p}(z)=\omega_i^{e_i}
       \exp\left(Cn_j\prod_{r=2}^d(z^r-1)
                                      \log(z/\omega_i)\right).
 \tag{A.66.21}
\]
Writing $z=\omega_i(1+pt)$, the logarithm of its principal
unit part, divided by $p$, is
\[
 Cn_j\prod_{r=2}^d(\omega_i^r-1)t+pG_i(t),
                     \qquad G_i\in\Z_p\langle t\rangle.
 \tag{A.66.22}
\]
The linear coefficient is a unit.  Thus this map is a
bijection from $\omega_i(1+p\Z_p)$ onto
$\omega_i^{e_i}(1+p\Z_p)$.

When $p\le B(D_j)$, choose a generator $\lambda$ of the
common power image in (A.66.19).  Every exponent map has
a primitive preimage of $\lambda$.  In detail, write
$N=p-1$, $g=\gcd(e_i,N)$ and $M=N/g$; solve
$(e_i/g)k=t\pmod M$ for the unit $t$ specifying
$\lambda$.  The reduction map
$(\Z/N)^\times\to(\Z/M)^\times$ is onto by the Chinese
remainder theorem, so choose $k$ a unit modulo $N$.
For $M=1$ any primitive element suffices.
The resulting $\omega_i$ are noncritical and
$F_{i,p}(\omega_i)=\lambda$.  The high error in
(A.66.14), followed by an adjustment in a unit coefficient
using (A.66.22), gives the exact equation (A.66.20), with
$F_{i,p}$ in its three coordinates.

When $p>B(D_j)$, take the existing unit $D_j$-power point
from the invariant.  There are at most
$M_d=\sum_{r=2}^d r$ critical nonzero residues for each
coordinate, and hence at most $M_d$ forbidden target
values under each exponent map.  A common scaling of
the three $D_j$-power values avoids all at most $3M_d$
targets.  Indeed their common image has size at least
$(p-1)/D_j>3M_d$, because $D_j\ge d$ and
\[
 \frac{(D+1)^4-1}{D}=D^3+4D^2+6D+4
                         >3\sum_{r=2}^d r\quad(D\ge d).
 \tag{A.66.23}
\]
Lift the scaling to an actual unit and scale the existing
local point.  Each of its three $D_j$-power values now has
a noncritical preimage under its exponent map; (A.66.22)
then gives an exact local input for that value.
This works also at an old coefficient prime: the invariant
requires the three local variables to be units, and does
not require the three coefficients to be units.

Approximate these local inputs modulo $p^{L_p}$.
On their indicated balls all the maps just used are
$1$-Lipschitz: this follows from their integral restricted
series, with the factor $A_p/p^s$ integral in (A.66.9).
The residues retain (A.66.16) and make the updated sum
error at least $H_{j+1}(p)$.  For an integer prime in
the final CRT class, (A.66.21) equals its actual integer
power $Q_i^{E_j(Q_i)}$: its root-of-unity part is fixed
by the earlier residues modulo $p-1$, and its principal
unit part is the displayed exponential.

\emph{Realizing the residues by primes.}
Put $n=n_j$, and use the three old coefficient fields
\[
 \begin{split}
 L_q&=\Q(\mu_n,(b_j/c_j)^{1/n}),\\
 L_r&=\Q(\mu_n,(a_j/c_j)^{1/n}),\\
 L_s&=\Q(\mu_n,(-a_j/b_j)^{1/n}).
 \end{split}
 \tag{A.66.24}
\]
The exact radical-field calculation in A.59 applies:
the coefficient shapes (A.66.13) give squareclasses
$-2,3,6$, and their distinguished valuations at two
or three are $1$ or $-1$.  Each field has maximal
abelian subfield exactly $\Q(\mu_n)$.  More explicitly,
its cyclic geometric kernel has order $n/2$; its
commutator with the cyclotomic group contains the
squares of that kernel, and the remaining quartic
field contains the nonabelian splitting field of a
fourth root of $2,3$, or $6$.  This forces the entire
kernel to be the commutator subgroup, as proved in A.59.

All three CRT classes are $1$ modulo $n$.
If $M=\prod_{p\le P_0}p^{L_p}$, it is divisible by $n$,
and the exact intersection is
$L_i\cap\Q(\mu_M)=\Q(\mu_n)$.
Thus the identity on $L_i$ and the desired CRT automorphism
of $\Q(\mu_M)$ agree on their intersection.  Chebotarev,
in precisely the compositum form used in A.59, supplies
infinitely many primes in that CRT class splitting in
$L_i$.  Choose three distinct such primes $q_j,r_j,s_j$
greater than $P_0$, and update
\[
 a_{j+1}=a_jq_j^{E_j(q_j)},\qquad
 b_{j+1}=b_jr_j^{E_j(r_j)},\qquad
 c_{j+1}=c_js_j^{E_j(s_j)}.
 \tag{A.66.25}
\]
The fields and congruences all used the old coefficients;
no congruence between the three new primes is required.

\emph{Checking every prime at the next stage.}
At each processed prime the new error is at least
$H_{j+1}(p)>2v_p(D_{j+1})$.  Vary a coordinate with
unit coefficient near $(1,1,1)$; the one-variable
Hensel inequality gives a unit $D_{j+1}$-power point.

For any of the three new primes $Q$, the imposed residues
give the exact identity
\[
                         \gcd(D_{j+1},Q-1)=n_j.
 \tag{A.66.26}
\]
At a prime factor of $n_j$ this follows from its prescribed
valuation in $Q-1$; every new prime divisor of $n_{j+1}$
was processed in (A.66.16) and has $Q\ne1$ modulo it.
Every prime divisor of $Cd!$ already divides $n_j$.
Consequently the $D_{j+1}$th and $n_j$th power subgroups
of $\F_Q^\times$ coincide.

At $q_j$, the unit ratio $c_{j+1}/b_{j+1}$ is an
$n_j$th power modulo $q_j$: this follows from splitting
in $L_q$ and the divisibility of every update exponent
by $n_j$.  Equation (A.66.26) makes it a $D_{j+1}$th
power.  Set $X=Z=1$ and lift a root for $Y$ by Hensel;
the term $a_{j+1}$ vanishes modulo $q_j$, and the
derivative is a unit.  At $r_j$ use $c_{j+1}/a_{j+1}$
and solve for $X$ with $Y=Z=1$.  At $s_j$ use
$-a_{j+1}/b_{j+1}$ and solve for $Y$ with $X=Z=1$.
All three primes exceed every prime factor of $D_{j+1}$.

At every remaining prime $p>P_0$ all coefficients are
units and $p>B(D_{j+1})$.  The finite-field character-sum
bound proved in A.44 gives a nonzero affine
$D_{j+1}$-power solution; its unit derivative lifts it
to $\Q_p$.  This covers every prime and proves the
invariant (A.66.13)--(A.66.14) at stage $j+1$.

\emph{The completed classes and ordinary local points.}
Every new exponent is divisible by $n_j$, which tends
to zero in $\widehat\Z$.  Thus the rational pairs
$(a_j/c_j,b_j/c_j)$ converge in the completed Kummer
group to (A.66.2).  At a fixed prime $p$ its valuation
in each of $a_j,b_j,c_j$ is eventually constant.
All later factors are units, and their factorial
exponents eventually annihilate every finite quotient
of $\Z_p^\times$.  The three products consequently
converge to nonzero numbers $a_p,b_p,c_p\in\Q_p$.
Eventually $p\le B(D_j)$, while $H_j(p)\to\infty$,
so (A.66.14) gives $a_p+b_p=c_p$.
Set $x_p=a_p/c_p$; then $1-x_p=b_p/c_p$ and both
are nonzero.  The local Kummer classes of these limits
are the restrictions of (A.66.2), since every local
power subgroup is open and the displayed ordinary
sequences converge in $\Q_p^\times$.

The signs of the global representatives are positive
and negative, respectively, giving the real classes
of $(3,-2)$.  Each of the infinitely many new primes
has exactly the valuation pair in (A.66.3), because
no later prime repeats it.  A rational point has
finitely many nonzero valuations.  The zero cusp has
$B=1$, contrary to $v_2(B)=1$; the one cusp has
$A=1$, contrary to $v_3(A)=1$; and the infinity cusp
has equal valuation vectors, contrary to the pair at
two.  Finally put $e_j(Q)=E_j(Q)/C$.  The products
\[
 A_0=\prod_{j\ge0}\frac{q_j^{e_j(q_j)}}{s_j^{e_j(s_j)}},
 \qquad
 B_0=\prod_{j\ge0}\frac{r_j^{e_j(r_j)}}{s_j^{e_j(s_j)}}
 \tag{A.66.27}
\]
converge in the completed group because every exponent
$e_j(Q)$ is still divisible by $n_j$.  They give the
actual compatible roots $A=3A_0^C$, $B=-2B_0^C$;
in particular reduction modulo $C$ is exactly $(3,-2)$.
\end{proof}
